
Benchmark contamination in foundation model training threatens evaluation validity. Recent work has shown that membership inference achieves near-random performance during pretraining (AUC≈50\%) and that adversarial paraphrasing evades semantic similarity detection (TPR<2\% at 1\% FPR) while preserving performance gains. We designed a three-tier detection architecture combining data-layer filters, Task Signature Graph probes, and geometric trajectory analysis. The system was intended to achieve ≥80\% detection power at <5\% false positive rate by analyzing learning dynamics rather than instance-level features. However, the experiment failed before testing theoretical claims: the detector flagged 100\% of samples (both clean and contaminated) because all three tiers returned constant positive values instead of computing detection metrics. Analysis revealed that 67\% of planned implementation tasks were not executed. The detector passed data validation checks (real datasets loaded correctly) but did not implement detection algorithms. This represents a validation scope gap: data correctness was verified independently of processing correctness. We document this failure to inform validation methodology for computational research. No theoretical advances in contamination detection resulted from this work.
